The analysis uses standard ATLAS reconstruction techniques to reconstruct leptons and jets, and to calculate the missing transverse momentum, enabling the identification of potential events of interest before employing a dedicated tracklet algorithm to specifically identify events containing disappearing tracks.

Standard tracks are reconstructed in the ID within $|\eta| < 2.5$~\cite{PERF-2015-08} and subsequently combined to create primary vertex candidates formed from at least two tracks with $\pt> 500\,\MeV$~\cite{ATL-PHYS-PUB-2015-026}. Among these, the primary vertex is identified as the vertex with the largest $\sum{\pt^2}$, where the sum runs over tracks associated with the vertex. All events considered in the analysis are required to have a primary vertex. The transverse and longitudinal impact parameters of all tracks, denoted by $d_0$ and $z_0$ respectively, are calculated relative to the beamline. The $z$-position of the primary vertex is used to calculate the point of closest approach, which is used as the reference point to calculate $z_0$.

Electron candidates are reconstructed from clustered energy deposits in the electromagnetic calorimeter matched to tracks re-fitted to account for bremsstrahlung losses~\cite{EGAM-2018-01}. They are required to have $\pt > 10\,\GeV$ and $|\eta| < 2.47$, and must satisfy the \enquote{LooseAndBLayer} quality criteria~\cite{PERF-2017-01}. In order to ensure that the trajectories of the electrons are consistent with originating from the primary vertex, the longitudinal impact parameter must satisfy $|z_0\sin\theta| < 0.5\,$mm, while the significance of the transverse impact parameter, $|d_0|/\sigma(d_0)$, must be less than 5. Electrons are required to be isolated from other objects by assessing a combination of track- and calorimeter-based information at the \enquote{FCLoose} working point~\cite{EGAM-2021-01}.

Muon candidates are reconstructed by combining ID tracks with matching MS tracks. They are required to have $\pt > 10$~\GeV and $|\eta| < 2.7$, and must satisfy the \enquote{Medium} quality criteria~\cite{MUON-2018-03}. Muons originating from the primary vertex are selected by requiring $|z_0\sin\theta| < 0.5\,$mm and $|d_0|/\sigma(d_0) < 3$. The muons must satisfy isolation requirements similar to those applied to electrons, but using the \enquote{Loose\_FixedRad} working point~\cite{MUON-2022-01}. Muons reconstructed purely from MS information are used to estimate muon backgrounds, and are referred to as stand-alone muons.

Hadronic jets are reconstructed using the \antikt algorithm~\cite{Cacciari:2008gp} as implemented in \texttt{FastJet}~\cite{Fastjet} with a jet radius parameter $R = 0.4$. The input to this algorithm consists of particle flow~(PFlow) objects~\cite{PERF-2015-09}, which combine measurements from the ID and calorimeters~\cite{PERF-2014-07} to improve the jet energy resolution and increase the jet reconstruction efficiency, especially at low jet $\pT$. Calibrations are applied to the jet energy scale and resolution, and these include components derived from both simulation and in situ measurements~\cite{JETM-2018-05}. The analysis uses jets within $|\eta| < 2.8$ with a calibrated $\pt > 20\,\GeV$. In order to reduce contributions from jets produced in \pileup interactions, all jets with $|\eta| < 2.5$ and $\pt < 60\,\GeV$ are required to satisfy the jet-to-vertex tagger~(JVT)~\cite{PERF-2014-03} \enquote{Tight} working point requirements. The tagger is configured to identify jets from the hard-scatter vertex with 92\% efficiency.

An overlap removal procedure is applied to avoid double counting of energy from different types of objects. If an electron and a jet are reconstructed within $\Delta R < 0.2$, the jet is rejected, while the electron candidate is kept. Subsequently, if an electron or muon is separated from a jet by $\Delta R < 0.4$, the electron or muon is discarded and the jet is kept.

The missing transverse momentum, with magnitude \met, is calculated as the negative vector sum of the transverse momenta of all identified hard physics objects (electrons, muons, jets), plus a contribution from an additional \enquote{soft term}. Forward jets within $2.8 < |\eta| < 4.5$ are also included in the \met calculation. The soft term is calculated from tracks matched to the primary vertex and not assigned to any of the hard objects~\cite{JETM-2020-03} but does not contain the reconstructed tracklets or charged pions described in the following paragraphs. 

To reconstruct the disappearing track due to a charged particle decaying after passing through the pixel detector, a dedicated tracklet reconstruction algorithm is applied after the standard reconstruction techniques described above are performed~\cite{ATL-PHYS-PUB-2019-011}. The tracking-detector hits already associated with reconstructed standard tracks are excluded from this tracklet-finding iteration. The algorithm attempts to identify tracklets having three or four hits in the pixel detector. The procedure attempts to extend the tracks into the SCT and TRT detectors and vetoes the tracklet if it can be associated with additional hits. Compared to the previous ATLAS analysis using the disappearing-track signature~\cite{SUSY-2018-19}, the veto does not consider an association if more SCT hits are only found multiple layers away from the tracklet. This reduces cases where tracklets would have been vetoed due to spurious SCT hits~\cite{ATL-PHYS-PUB-2019-011} further into the detector. To improve the momentum resolution of the tracklets, they are re-fitted, adding the primary vertex as an additional constraint, and the tracklet is required to have $|d_0|/\sigma(d_0) < 1.5$ and $|z_0\sin\theta|< 0.2$\,mm. An overlap removal procedure discards tracklets within $\Delta R = 0.4$ of a jet, a lepton, or a track in the MS. Isolation requirements are also applied to the tracklets: the scalar sum of the \pt of all tracks within a cone of radius $\Delta R=0.4$ around the tracklet, $p_{\mathrm{T}}^{\mathrm{cone40}}$, must satisfy $p_{\mathrm{T}}^{\mathrm{cone40}}/p_{\mathrm{T}}^{\mathrm{tracklet}} < 0.04$ and the total transverse energy of all topological clusters within the $\Delta R=0.4$ cone, $E_{\mathrm{T}}^{\mathrm{topoclus40}}$, must be less than 5\,\GeV. Tracklets are required to have $|\eta| > 0.1$, vetoing the central region where MS coverage is reduced. Figure~\ref{fig:reco:TrackEff} shows that the reconstruction efficiency for three- and four-layer tracklets can reach 80\% depending upon the distance travelled by the chargino before decaying. 

To identify low-energy charged pions, indicative of the electroweakino signal scenario, a boosted decision tree (BDT) machine-learning algorithm is used. After performing the tracklet reconstruction described above, a set of low-momentum tracks (\enquote{soft tracks}) are reconstructed from the remaining hits. These low-momentum tracks are seeded from at least three SCT hits that are within $\Delta R = 0.8$ of a tracklet and are then extended using further SCT or TRT hits. A minimum of six hits, and at most two holes\footnote{A hole is a missing hit on a reconstructed track trajectory, possibly due to a non-functioning strip which was not masked.}~\cite{ATL-PHYS-PUB-2019-011}, are required to reconstruct a soft track and a requirement of $\pt > 200$\,MeV is applied. As only a limited number of four-layer tracklet events have a matching soft-track that satisfies the soft-track reconstruction requirements, only events containing tracklets produced with three pixel hits are used when performing the pion reconstruction. The BDT is used to select viable charged pion candidates and to reject events with coincidental hits forming a soft track. The algorithm is trained with MC signal simulation, using soft tracks that have been \enquote{truth-matched} at generator level to charged pions from long-lived chargino decays using a $\Delta R$-matching criterion.

The background is taken from data, selecting tracklets that have large impact parameter values with respect to the primary vertex. From MC studies this is expected to be dominated by accidental crossings, referred to as the \enquote{fake track} background. The algorithm uses eight input variables related to the properties of the soft tracks, including the kinematic properties \pt and $\eta$, the number of hits/holes associated with the track in the SCT, the number of hits associated with the track in the TRT, the distance in $\Delta R$ between the soft track and the tracklet, the distance in $z$ between the soft track and the tracklet, and the $\chi^2$ value of the soft track when fitted to the tracklet's endpoint. The TMVA architecture~\cite{TMVA} is used with gradient boost, and a hyperparameter optimisation is performed to prevent overtraining. Figure~\ref{fig:reco:pionBDT} displays the BDT output score. A low-energy charged pion is identified if the score for a soft track is ${>}0.8$ and, to reduce the misidentification rate, a transverse impact parameter requirement of $|d_0|/\sigma(d_0) < 40$ is satisfied. The latter rejects an additional 4\% of the tracklet background with minimal impact on the signal, according to MC simulation. Figure~\ref{fig:reco:pionEff} shows that the pion reconstruction efficiency, measured as a function of the true pion \pt in the lab-frame, is as high as 70\% for \pt near 2\,\GeV.

\begin{figure}[tb!]
\centering
  \begin{subfigure}{0.44\textwidth}
  \includegraphics[width=\textwidth]{../figures/reco/Public/decayr_250_0p04.pdf}
  \caption{}
    \label{fig:reco:TrackEff}
  \end{subfigure}
  \begin{subfigure}{0.44\textwidth}
  \includegraphics[width=\textwidth]{../figures/reco/Public/bdt_score_3layer_scteffective.pdf}
      \label{fig:reco:pionBDT}
  \end{subfigure}
  \begin{subfigure}{0.44\textwidth}
  \includegraphics[width=\textwidth]{../figures/reco/Public/sp_eff_250_0p04.pdf}
      \label{fig:reco:pionEff}
  \end{subfigure}

\caption{Representative distributions related to the reconstruction of tracklets and charged pions: \protect\subref{fig:reco:TrackEff} event yields and reconstruction efficiencies for  three- and four-layer tracklet reconstruction as a function of \chinoonepm decay radius; \protect\subref{fig:reco:pionBDT} BDT output score used to judge if an event contains a low-energy pion, showing the training dataset (solid histogram) and the test dataset (points); \protect\subref{fig:reco:pionEff} event yields and reconstruction efficiency for charged pions as a function of the true pion transverse momentum as measured in the lab-frame, where the acceptance includes all charged-pion reconstruction effects but not geometric acceptance from tracklet reconstruction (which is assumed to already have been reconstructed). The true transverse momentum is defined as $\pt = \sqrt{p_x^2 + p_y^2}$ in the global detector coordinate system with the $z$-axis along the beam pipe. For \protect\subref{fig:reco:TrackEff} and \protect\subref{fig:reco:pionEff} the filled histogram displays the total yield without applying a selection, while the dashed lines present the yield after applying the tracklet selection or tracklet+pion selection. The points denote the reconstruction efficiency. A representative higgsino signal scenario with $m(\chinoonepm) = 250$\,GeV and $\tau(\chinoonepm)= 0.04$\,ns is shown. 
 }
\end{figure}

% taken from INT table 18
% Just in case we want to be explicit here are all of the things
%The algorithm uses 8 input variables: $\eta^{\mathrm{soft-trk}}$, $p_{\mathrm{T}}^{\mathrm{soft-trk}}$, $N_{\mathrm{SCT-hit}}^{\mathrm{soft-trk}}-N_{\mathrm{SCT-Shared-Hit}}^{\mathrm{soft-trk}}$, $N_{\mathrm{SCT-holes}}^{\mathrm{soft-trk}}$, $N_{\mathrm{TRT-hit}}^{\mathrm{soft-trk}}$, the $\chi^2_{\mathrm{secondary-vertex}}$ fit value, the longitudinal distance in the transverse plane between the crossing of the soft track and the tracklet ($d_{z}^{x-y}$) and the distance between the tracklet and soft-track ($\Delta R$). 



\FloatBarrier
